\subsection{R-27: exact four-variable dimension from a uniform endpoint repair}

\paragraph{Statement R-27.}
For every $m\geq17$ there are a ternary Jacobi pencil $A$ and
linear vectors $u,v$, supported at its two endpoints, for which
\begin{equation}\label{R27:criterion}
 \dim\mathcal C_A(u,v)=8m+2.
\end{equation}
Consequently, with the notation of R-13 and R-23,
\begin{equation}\label{R27:dimension}
 c_m=\min\left\{\binom{m+3}{3},\,4m^2-2m+2\right\}
 \quad(m\geq1),
 \qquad \epsilon_m=0\quad(m\geq17).
\end{equation}
The uniform assertion follows from one fixed finite boundary
certificate and an insertion argument. No extrapolation of
finite-size ranks, and no assertion about variable counts growing
with $m$, is used.

\begin{proof}
We use the notation and continuants of R-26. In particular,
$L=m-1$, $h=yz$, $A_{ii}=\alpha_i x$, $A_{i,i+1}=y$,
$A_{i+1,i}=\delta_i z$, and
\[
 D_i=\prod_{e=i}^{L-1}\delta_e,\qquad
 E_i=\prod_{e=0}^{i-1}\delta_e,\qquad D=D_0=E_L.
\]
All parameters are initially nonzero. We work over
characteristic zero; reduction modulo $101$ will certify specified
nonzero rational minors.

\paragraph{1. General endpoint formulas and the seven base relations.}
Allow arbitrary endpoint linear forms
\[
 u=(U_0,0,\ldots,0,U_1)^t,\qquad
 v=(V_0,0,\ldots,0,V_1)^t.
\]
The path decomposition used in R-26 gives
\begin{align}
 a_i&=V_0y^iP_{[i+1,L]}+V_1D_i z^{L-i}P_{[0,i-1]},
                                                    \label{R27:a}\\
 b_i&=U_0E_i z^iP_{[i+1,L]}+U_1y^{L-i}P_{[0,i-1]}.
                                                    \label{R27:b}
\end{align}
For $0<i<L$ the second auxiliary derivatives are
\begin{align}
 t_i&=2V_0V_1\bigl(U_0D_i z^{L-i}P_{[1,i-1]}
                         +U_1y^iP_{[i+1,L-1]}\bigr),
                                                    \label{R27:t}\\
 t'_i&=2U_0U_1\bigl(V_0y^{L-i}P_{[1,i-1]}
                         +V_1E_i z^iP_{[i+1,L-1]}\bigr).
                                                    \label{R27:tp}
\end{align}
At $i=0$ omit the first summand in each parenthesis, and at
$i=L$ omit the second. These are exact polynomial identities:
each complementary matching has a forced path and an independent
interval matching, while the two assignments of the rank-one
entries give the factor $2$.

Let $C(\varepsilon)$ be the coefficient matrix of the $8m+3$
columns defining $\mathcal C_A(u,v)$, when
\begin{equation}\label{R27:endpoint_line}
 (U_0,U_1,V_0,V_1)
 =(x,y,x,z)+\varepsilon(y+2z,\,3x+4z,\,5y+6z,\,7x+8y).
\end{equation}
Its entries have degree at most three in $\varepsilon$. Write
$C=C_0+\varepsilon C_1+\varepsilon^2C_2+\cdots$.
At $\varepsilon=0$ the character decomposition of R-26 applies.
There are eight columns in each character $2,\ldots,m-2$ and
nine columns in each of $0,1,-1$.

For later upper bounds we give all seven independent base
relations. Three have character zero:
\begin{align}
 xa_0+ya_L-xb_0-zb_L&=0,\label{R27:scaling_base}\\
 K_{00}:=yb_1-\delta_0za_1+\alpha_0xa_0-\alpha_0xb_0&=0,
                                                    \label{R27:K00}\\
 K_{01}:=t_0-t'_0&=0.                              \label{R27:K01}
\end{align}
Two have character one:
\begin{align}
 K_{10}&:=\delta_0t_1-2xb_L+\alpha_0t'_L=0,
                                                    \label{R27:K10}\\
 K_{11}&:=\alpha_L\delta_0
       (\delta_1za_2-ya_0+\alpha_1xa_1)\notag\\
 &\hspace{17mm}-\alpha_0\alpha_1
       (\delta_{L-1}zb_{L-1}-yP+\alpha_Lxb_L)=0.
                                                    \label{R27:K11}
\end{align}
The remaining two are the transposed and gauged versions of
\eqref{R27:K10}--\eqref{R27:K11}, of character $-1$.
Here and below identities involving these $K$'s without a stated
endpoint specialization mean the displayed constant linear
combination of the corresponding columns.

For precision, the symmetry producing the two negative-character
relations is as follows. Transpose $A$, exchange $y,z$, and
conjugate by the diagonal matrix with entries $1/E_i$, in the
form $S^{-1}AS$. The upper and lower coefficients again become
$1,\delta_i$, while the endpoint pair becomes
$(U_0,U_1,V_0,V_1)=(x,Dy,x,D^{-1}z)$.
Next make the variable change $(y,z)\mapsto(ty,t^{-1}z)$ and
conjugate by $\operatorname{diag}(t^{-i})$, choosing
$t^m=D^{-1}$. This restores $(x,y,x,z)$. These operations are
invertible over $\mathbb C$, preserve the Jacobi family, exchange
characters $r,-r$, and take arbitrary endpoint perturbations to
arbitrary endpoint perturbations.

All five displayed identities follow by the left or right
continuant recurrence. For example, the two parentheses in
\eqref{R27:K11}, at the base endpoints, equal respectively
$2\alpha_0\alpha_1D_1x^2z^{m-1}$ and
$2\alpha_LDx^2z^{m-1}$. The three character-zero relations are
independent: the Hessian columns distinguish $K_{01}$, and the
columns with index one distinguish $K_{00}$ from scaling.
In character one, $K_{10}$ has Hessian support and $K_{11}$ has
a nonzero $yP$ coefficient. The same argument applies after the
symmetry. Thus the seven relations are independent whenever the
parameters are nonzero. R-26 proves the reverse generic rank
bound. Consequently
\begin{equation}\label{R27:base_rank}
 \operatorname{rank}C_0=8m-4
\end{equation}
on a nonempty open set, and its kernel consists of these seven
relations, entirely in characters $0,1,-1$.

\paragraph{2. A universal upper bound of four for the first normal.}
This step is needed to control the second-order continuation.
It holds for arbitrary linear endpoint perturbations, not just
\eqref{R27:endpoint_line}.
For general endpoint forms, direct use of the same recurrences
gives the exact identity
\begin{equation}\label{R27:K11_general}
 K_{11}=2\alpha_0\alpha_1\alpha_LD\,
       xz^{L-1}(xV_1-zU_0)
          -\alpha_0\alpha_1(U_1-y)P.
\end{equation}
Indeed, its first parenthesis is
$2\alpha_0\alpha_1D_1x^2z^{L-1}V_1$, and its second is
$2\alpha_LDxz^LU_0+(U_1-y)P$. This proves
\eqref{R27:K11_general}, including every endpoint factor.

We show that the first derivative of this identity lies in
$\operatorname{im}C_0$. The term involving $P$ already does.
The other terms are linear combinations of
\begin{equation}\label{R27:five_pure}
 x^3z^{m-2},\quad x^2yz^{m-2},\quad x^2z^{m-1},
 \quad xyz^{m-1},\quad xz^m.
\end{equation}
Here are explicit base-column witnesses for all five.
For $r=2,3$, put $j=m-r$ and define
\begin{align}
 R_A(r)&=ya_{r-1}-\alpha_rx a_r-\delta_rza_{r+1}
       =-2\alpha_rD_r xz^{m-r}P_{[0,r-1]},
                                                    \label{R27:RA}\\
 R_B(r)&=yb_{j+1}-\alpha_jxb_j-\delta_{j-1}zb_{j-1}
       =-2\alpha_jE_jx^2z^jP_{[j+1,L]}.
                                                    \label{R27:RB}
\end{align}
The unwanted path branch cancels by one continuant recurrence;
the remaining branch is the displayed expression by the recurrence
at the other end. When $r=2$, the two residual polynomials span
$x^3$ and $xh$, because
$P_{[0,1]}=\alpha_0\alpha_1x^2+\delta_0h$ and
$P_{[L,L]}=\alpha_Lx$. Hence the first and fourth forms in
\eqref{R27:five_pure} are available. When $r=3$ they span
$x^4$ and $x^2h$ provided
\begin{equation}\label{R27:Delta}
 \Delta=\alpha_0\alpha_1\alpha_2\delta_{L-1}
       -\alpha_{L-1}\alpha_L(\alpha_2\delta_0+\alpha_0\delta_1)
       \ne0.
\end{equation}
This follows by expanding $xP_{[0,2]}$ and
$x^2P_{[L-1,L]}$; their coefficient determinant is precisely
$\Delta$. It supplies the second form in \eqref{R27:five_pure}.
Finally
\begin{align*}
 \delta_1za_2-ya_0+\alpha_1xa_1
       &=2\alpha_0\alpha_1D_1x^2z^{m-1},\\
 \delta_0za_1-xP+\alpha_0xa_0
       &=2\alpha_0D xz^m
\end{align*}
supply the other two. Condition \eqref{R27:Delta} defines a
nonempty open set: when all parameters equal one, $\Delta=-1$.

Thus $K_{11}$ has zero first normal on a nonempty open set.
The symmetry above proves the same assertion for its independent
negative-character counterpart. Scaling also has zero first
normal, since the actual polynomial relation
\begin{equation}\label{R27:scaling_actual}
 U_0a_0+U_1a_L-V_0b_0-V_1b_L=0
\end{equation}
holds for every endpoint tuple; differentiating it gives its
first-order lift. These are three independent vectors in the
seven-dimensional base kernel. Therefore, generically,
\begin{equation}\label{R27:first_upper}
 \operatorname{rank}\left(\ker C_0\longrightarrow
         \operatorname{coker}C_0,\quad k\longmapsto[C_1k]\right)
 \leq4.
\end{equation}
The witnesses above for $K_{11}$ use only the complete character
blocks $0,1,2,3$, and the transposed witnesses use only
$0,-1,-2,-3$.

\paragraph{3. The finite-order rank lemma and the complete boundary projection.}
We record the elementary linear algebra used below. Let
$M(\varepsilon)=M_0+\varepsilon M_1+\varepsilon^2M_2+\cdots$.
Let $r_0=\operatorname{rank}M_0$ and let $r_1$ be the rank of
\[
 N_1:\ker M_0\longrightarrow\operatorname{coker}M_0,
 \qquad k_0\longmapsto[M_1k_0].
\]
For $k_0\in\ker N_1$ choose $k_1$ with
$M_0k_1=-M_1k_0$. Define
\[
 N_2(k_0)=[M_2k_0+M_1k_1]
 \quad\hbox{in }\operatorname{coker}M_0/\operatorname{im}N_1,
 \qquad r_2=\operatorname{rank}N_2.
\]
Changing the choice of $k_1$ changes this class by an element
of $\operatorname{im}N_1$, so the map is well defined.
For the lower triangular coefficient matrices
\begin{equation}\label{R27:Toeplitz}
 T_0=M_0,\quad
 T_1=\begin{pmatrix}M_0&0\\M_1&M_0\end{pmatrix},\quad
 T_2=\begin{pmatrix}M_0&0&0\\M_1&M_0&0\\M_2&M_1&M_0\end{pmatrix}
\end{equation}
one has
\begin{equation}\label{R27:jet_ranks}
 \operatorname{rank}T_0=r_0,\quad
 \operatorname{rank}T_1=2r_0+r_1,\quad
 \operatorname{rank}T_2=3r_0+2r_1+r_2.
\end{equation}
Moreover, $\operatorname{rank}_{\mathbb C(\varepsilon)}M
\geq r_0+r_1+r_2$. To prove these assertions, choose constant
bases that put $M_0$ in an identity block with a zero complement.
Eliminate that block by invertible formal row and column
operations. The complement starts with $\varepsilon N_1$.
Eliminate an identity block of $N_1$ in the same way; the
remaining coefficient of $\varepsilon^2$ is exactly the map
$N_2$ above. On truncation modulo $\varepsilon^{s+1}$,
invertible formal changes give invertible triangular block
changes. A unit pivot contributes $s+1$ to the rank, a pivot
of order one contributes $s$, and a pivot of order two contributes
$s-1$ when $s\geq2$. This proves all the formulas and the
generic-rank lower bound without an analytic limiting argument.

Now retain every base column of $C$ whose character is in
\[
 B_m=\{0,\pm1,\ldots,\pm5\}\pmod m,
\]
and retain every monomial row of these same characters. There
are $91$ retained columns: nine in each of $0,1,-1$, and eight
in each of the other eight characters. Denote this projected
matrix by $S(\varepsilon)$. On the generic base-rank open,
$\operatorname{rank}S_0=84$ and its kernel is precisely the
seven-dimensional full kernel.

It is essential here to retain complete character blocks.
Replacing one endpoint monomial by another linear monomial
changes a character by at most two. Consequently, for every
$k_0\in\ker C_0$, $C_1k_0$ is supported in characters
$0,\pm1,\pm2,\pm3$. If its class modulo the full base image
vanishes, its correction $k_1$ can be chosen entirely in these
complete base character blocks: project any solution of
$C_0k_1=-C_1k_0$ onto them. Then
\begin{equation}\label{R27:support}
 C_2k_0+C_1k_1
 \quad\hbox{is supported in characters }0,\pm1,\ldots,\pm5.
\end{equation}
Indeed, two replacements shift the original kernel characters
by at most four, and one replacement shifts the correction
characters by at most two. The first and second normal maps
of $S$ are therefore exactly the first and second normal maps
of the full $C$, with their images in these subspaces. Columns
of outside base characters cannot change either quotient.
This also shows why a rank assertion for a chosen boundary
submatrix alone would not suffice.

\paragraph{4. A fully specified size17 boundary certificate.}
Use $m=17$, $L=16$, and the following integer parameters:
\begin{align*}
 (\alpha_0,\ldots,\alpha_{16})
   &=(86,18,1,17,95,90,80,90,48,100,9,60,2,25,78,53,54),\\
 (\delta_0,\ldots,\delta_{15})
   &=(34,1,37,24,47,1,59,41,0,92,38,56,75,6,44,14).
\end{align*}
The zero is the central edge $8$. The endpoint line is exactly
\eqref{R27:endpoint_line}. The calculation at a zero edge is
interpreted by the following rational row normalization, rather
than as the rank of the literal degenerate pencil.

For nonzero edges divide each retained row of integer weight
$w=b-c$ with $-m-5\leq w\leq-m+5$ by $D=\prod\delta_i$.
Leave every other row unchanged. Then take the limit
$\delta_8\to0$. In these negative extreme rows every term
contains the long lower path across edge $8$, and its residual
interval avoids the edge. The factor $\delta_8$ cancels before
taking the limit. Equivalently, for these rows set
$\delta_8=1$ and divide by the product of the other edges;
for the other rows set $\delta_8=0$. This product is nonzero
and is $1$ modulo $101$. This rule specifies every normalized
coefficient without division by zero.

Order degree18 monomials $x^ay^bz^c$ by decreasing $a$ and
then decreasing $b$, and retain the $124$ rows with
$(b-c)\bmod17\in B_{17}$. Order the original columns as
\[
\begin{aligned}
 &xP,yP,zP;\qquad xa_i,ya_i,za_i\quad(i=0,\ldots,16);\\
 &xb_i,yb_i,zb_i\quad(i=0,\ldots,16);\\
 &t_i\quad(i=0,\ldots,16);\qquad t'_i\quad(i=0,\ldots,16).
\end{aligned}
\]
and retain the $91$ columns of base character in $B_{17}$.
Their characters are, respectively,
$(0,1,-1)$, $(i,i+1,i-1)$, $(-i,-i+1,-i-1)$, $i$, and $-i$.
Expand the exact formulas \eqref{R27:a}--\eqref{R27:tp} in
$\varepsilon$, and form the three matrices
\eqref{R27:Toeplitz} with $M_i=S_i$.

The following is a finite rational-minor certificate. All
arithmetic in this table is in $\mathbb F_{101}$:
\begin{center}
\begin{tabular}{rrrr}
matrix&dimensions&rank&selected determinant\\ \hline
$T_0$&$124\times91$&84&99\\
$T_1$&$248\times182$&172&74\\
$T_2$&$372\times273$&262&45
\end{tabular}
\end{center}
For an unambiguous specification of the selected minors, use
the following deterministic elimination rule. Scan columns
from left to right. In each column choose the first nonzero
entry among unused rows, swap that row into the next pivot
position, normalize its pivot to one, and eliminate entries
below it. Record the original row names in pivot order and the
selected column names in scan order. The selected determinant
in this ordered-row convention is the product of the pivots
before normalization. The remaining rows are zero at the end.
Together with the integer input, the continuant recurrence,
the endpoint formulas, and the row and column ordering, this
algorithm specifies the entire finite certificate. In particular
the table is not dependent on an unspecified numerical rank
tolerance or an unspecified choice of a minor.

The complete coefficient arrays and selected row and column
lists are also preserved in
\texttt{jacobi\_boundary\_repair\_m17.npz} and
\texttt{boundary\_toeplitz\_certificates.json}, with builders
\texttt{check\_jacobi\_boundary\_repair.py} and
\texttt{check\_boundary\_toeplitz.py}. A separate checker
reconstructed all $33{,}852$ entries of $S_0,S_1,S_2$ by direct
enumeration of path matchings and multiplication of polynomials
in $x,y,z,\varepsilon$, and independently obtained the three
ranks $84,172,262$. This is an actual finite coefficient and
rank check; the uniform argument is supplied next. Each listed
nonzero residue proves a nonzero rational minor in
characteristic zero.

\paragraph{5. Inserting a central vertex preserves every retained jet.}
Start with the preceding central-cut data and insert an isolated
vertex with diagonal $x$ at edge $8$, replacing that zero edge
by two zero edges. The size increases by one. The left endpoint
indices are unchanged; the right endpoint indices increase by
one. Repeat at the leftmost central cut as often as desired.
All other coefficients retain their positions relative to the
appropriate endpoint. This constructs cut data in every size
$m\geq17$.

We verify the effect on the entire normalized boundary matrix
for arbitrary endpoint forms, so in particular on each of its
$\varepsilon$ coefficients. Retained first-derivative columns
have index at most $6$ from the left or at least $m-6$ from the
left. Retained Hessian columns have index at most $5$ or at
least $m-5$. In formulas \eqref{R27:a}--\eqref{R27:tp}, each
such column is the sum of a short-path term with a long residual
continuant and a long-path term with a short residual continuant.
Every long residual continuant contains the new vertex; every
short residual continuant avoids it. Every long forced path
crosses the central cut; every short forced path avoids it.
These assertions include the intervals $[i+1,L-1]$ and
$[1,i-1]$: their short endpoint distances are at most $5$,
whereas the original cut is between vertices $8$ and $9$.
They remain true after every insertion.

There is no ambiguity in the retained rows about which term
occurs. A first-derivative column has an auxiliary factor of
degree two, so its short-path terms have integer weight of
absolute value at most $6+2=8$, whereas its long-path terms
have absolute weight at least $(m-7)-2=m-9$.
A Hessian column has an auxiliary factor of degree three;
the corresponding bounds are $5+3=8$ and $(m-6)-3=m-9$.
The columns $xP,yP,zP$ have integer weight $0,1,-1$.
The retained monomial rows have precisely the three disjoint
weight windows
\begin{equation}\label{R27:windows}
 [-5,5],\qquad [m-5,m+1],\qquad [-m-1,-m+5].
\end{equation}
For $m\geq17$, the central window cannot contain a long-path
term, and either extreme window cannot contain a short-path
term. Possible overlap of the short and long weight bounds
outside these windows is discarded by the projection.

It follows that the new normalized boundary coefficient matrix
is obtained from the old one by the following three injective
monomial maps on its disjoint row windows:
\begin{equation}\label{R27:row_insertion}
 \hbox{central: }f\mapsto xf,\qquad
 \hbox{positive extreme: }f\mapsto yf,\qquad
 \hbox{negative extreme: }f\mapsto zf.
\end{equation}
For the central terms this is the continuant recurrence at an
isolated vertex. A long upper path gains one factor $y$.
A long lower path gains one factor $z$ and the new middle edge
factors; division by $D$ cancels all those factors. The quotient
depends only on nonzero outer edge products, which are unchanged.
Thus the assertion also holds when there are several central
zero edges. New monomial rows outside the images of the three
maps have zero coefficient. The column list is naturally
reindexed, and still has $91$ members.

The three maps have distinct integer-weight ranges and are
injective within those ranges. Hence the normalized matrices
$S_i$ and all three block matrices $T_i$ retain their ranks
under every insertion. The finite certificate therefore gives,
at a normalized cut limit in every size,
\begin{equation}\label{R27:limit_ranks}
 \operatorname{rank}T_0\geq84,\qquad
 \operatorname{rank}T_1\geq172,\qquad
 \operatorname{rank}T_2\geq262.
\end{equation}
To interpret this as a nonzero-edge witness, replace each
central zero edge by $\tau$ times a nonzero constant. The
negative extreme row quotients have already canceled every
middle edge factor, and their remaining denominators are
nonzero outer products. All normalized entries therefore have
the indicated finite limit as $\tau\to0$. Each nonzero limit
minor is a nonzero rational minor for nonzero edges. The rank
conditions \eqref{R27:limit_ranks} consequently hold on nonempty
open subsets of the actual Jacobi parameter space. Normalizing
rows by a nonzero $D$ is invertible and does not alter their
rank or their first and second normal maps.

\paragraph{6. Combining the six directions at one actual pencil.}
Fix a size $m\geq17$. Intersect the three nonempty minor opens
from \eqref{R27:limit_ranks} with the base-rank open of R-26
and with \eqref{R27:Delta}. The parameter space of nonzero
$\alpha_i,\delta_i$ is irreducible, so this finite intersection
is nonempty. At one common actual pencil in this intersection,
the projected matrix has $r_0=84$ and, by
\eqref{R27:first_upper}, $r_1\leq4$. The $T_1$ lower bound
and \eqref{R27:jet_ranks} force $r_1=4$.
The $T_2$ lower bound then gives
\[
 262\leq3\cdot84+2\cdot4+r_2,
 \qquad\hbox{hence }r_2\geq2.
\]
The support argument \eqref{R27:support} identifies these with
four first and at least two second normal classes for the full
matrix. Together with \eqref{R27:base_rank}, the finite-order
rank lemma therefore gives
\[
 \operatorname{rank}_{\mathbb C(\varepsilon)}C(\varepsilon)
       \geq(8m-4)+4+2=8m+2.
\]
Relation \eqref{R27:scaling_actual} gives the reverse bound,
since there are $8m+3$ columns and its coefficient vector is
nonzero. A nonzero maximal minor is a polynomial in
$\varepsilon$, so some complex value of $\varepsilon$ attains
this rank. This proves \eqref{R27:criterion}.

The upper bound on the first normal is indispensable in this
argument. Merely knowing three lower rank bounds at a cut
would not preserve a claimed second-order gain under parameter
variation: a pair of order-two pivots can become pivots of
orders one and three. The independent algebraic bound
\eqref{R27:first_upper} rules out that possibility here.

Finally apply the exact bordered induction of R-25 to
\eqref{R27:criterion}:
\[
 c_{m+1}\geq c_m+8m+2\qquad(m\geq17).
\]
The torus upper bound $4m^2-2m+2$ has precisely the same
increment. Its equality at $m=17$ was certified in R-24,
so induction proves equality for every $m\geq17$.
R-23 identifies the resulting defect with the specified
degree-$m$ first Tor group, proving $\epsilon_m=0$ in that
range. R-13 gives $c_m=\binom{m+3}{3}$ for $1\leq m\leq16$.
These statements, and the two upper bounds, yield exactly
\eqref{R27:dimension}.
\end{proof}
